\begin{textsample}{NEG Dim 4 – human – Score -1.00 – es/es\_ef\_anos\_iniciais\_arte\_2\_p1.txt}  \label{ex:f4_neg_016}
Para a seleção e a organização do ensino de Arte, foram estabelecidos critérios que, no conjunto, possibilitam que os educandos compreendam as relações entre tempos e contextos sociais dos sujeitos na sua interação com a arte e a cultura.
De acordo com a Base \textbf{Nacional} Comum \textbf{Curricular} ( BNCC ), a prática investigativa constitui o modo de produção e organização dos conhecimentos em Arte. É no trajeto do fazer artístico que os alunos criam, experimentam, desenvolvem e percebem seus percursos poéticos e estéticos.
Os conhecimentos, processos e técnicas produzidos e acumulados ao longo do tempo em Artes Visuais, Dança, Música e Teatro contribuem para a contextualização dos saberes e das práticas artísticas ( BRASIL, 2017 ).

% matched lemmas: curricular, nacional
\end{textsample}
